% ---- pg-368 （书页 356） ----

若总体为多元正态分布, 在习题 8.9 中给出一个关于所有变量独立的大样本检验(即检验 $\boldsymbol{\Sigma}$ 的所有非对角元为 0).

\section{用主成分监控质量}  % 原书 8.6
\phantomsection\label{sec:8.6}

在 5.6 节, 我们介绍了多元控制图, 包括质量椭圆图和 $T^2$ 图. 今天, 用电子计算机和其他方法自动收集数据, 收集 10 或 20 个过程变量的数据, 已非罕见之事. 大的化学和药品公司据说测量超过 100 个过程变量, 它们包括在生产过程的不同场合下的温度、压力及重量等. 即使使用 10 个变量作监控, 也可搭配出 45 对变量生成质量椭圆. 显然, 要求有新的手段, 既能形象化显示重要的变量又依旧能灵敏检测变异的特别原因.

\mvsub{检查一组已知测量值的稳定性}

令 $\mathbf{X}_1,\mathbf{X}_2,\cdots,\mathbf{X}_n$ 是从均值 $\boldsymbol{\mu}$、协方差矩阵 $\boldsymbol{\Sigma}$ 的正态分布中抽取的随机样本. 我们考虑前两个样本主成分, $\hat{y}_{j1}=\hat{\mathbf{e}}_1'(\mathbf{x}_j-\bar{\mathbf{x}})$ 和 $\hat{y}_{j2}=\hat{\mathbf{e}}_2'(\mathbf{x}_j-\bar{\mathbf{x}})$. 可以考虑添加主成分, 只是两个样本主成分更容易形象地检查, 并且在任何两个成分中, 前两个都解释样本总方差最大的累积比.

如果一个过程在全部时间里都是稳定的, 只由于共同原因所引起的变化影响所测量的特征, 那么在这种情况下, 前两个主成分的值应该是稳定的. 反之, 如果主成分保持全时稳定, 那么影响该过程的共同因素有希望保持为常数. 为了用主成分监控质量, 我们分两步考虑过程. 第一步是构造一个数值对 $(\hat{y}_{j1},\hat{y}_{j2})$ $(j=1,2,\cdots,n)$ 的椭圆形式图.

根据\hyperref[eq:8-20]{式(8-20)}, 第一主成分 $\hat{y}_1$ 的样本方差由最大特征值 $\hat{\lambda}_1$ 给出, 第二主成分 $\hat{y}_2$ 的样本方差由第二大特征值 $\hat{\lambda}_2$ 给出. 两个样本主成分是不相关的, 故对大的 $n$, 质量椭圆(见 5.6 节)为满足下述条件的可能的数值对 $(\hat{y}_1,\hat{y}_2)$ 的点集:
\phantomsection\label{eq:8-36}
\begin{equation}
\frac{\hat{y}_1^2}{\hat{\lambda}_1}+\frac{\hat{y}_2^2}{\hat{\lambda}_2}\leqslant\chi_2^2(\alpha) \tag{8-36}
\end{equation}

\noindent\textbf{例 8.10}\quad (基于前两个主成分的椭圆形式图)

参看\hyperref[tab:5.8]{表 5.8} 给出的警署加班数据. \hyperref[tab:8.1]{表 8.1} 包含样本协方差矩阵 $\mathbf{S}$ 的五个规范化了的特征向量和特征值.

前两个样本成分解释总方差的 82\%.

\begin{mvtable}
\phantomsection\label{tab:8.1}\mvtabcap{表 8.1\quad 警署数据的样本协方差矩阵的特征向量和特征值}
\begin{tabular}{lccccc}
\toprule
变量 & $\hat{\mathbf{e}}_1$ & $\hat{\mathbf{e}}_2$ & $\hat{\mathbf{e}}_3$ & $\hat{\mathbf{e}}_4$ & $\hat{\mathbf{e}}_5$\\
\midrule
法定加班出勤时数 $(x_1)$ & 0.046 & $-0.048$ & 0.629 & $-0.643$ & 0.432\\
特殊事件处理加班时数 $(x_2)$ & 0.039 & 0.985 & $-0.077$ & $-0.151$ & $-0.007$\\
处理遗留人员加班时数 $(x_3)$ & $-0.658$ & 0.107 & 0.582 & 0.250 & $-0.392$\\
允许补偿的加班时数 $(x_4)$ & 0.734 & 0.069 & 0.503 & 0.397 & $-0.213$\\
会议时数 $(x_5)$ & $-0.155$ & 0.107 & 0.081 & 0.586 & 0.784\\
\midrule
$\hat{\lambda}_i$ & 2 770 226 & 1 429 206 & 628 129 & 221 138 & 99 824\\
\bottomrule
\end{tabular}
\end{mvtable}

\hyperref[tab:8.2]{表 8.2} 表示全部五个成分的样本值.

我们用前两个样本主成分构造 95\% 椭圆形式图, 并画出 16 对成分值的图形.
